% DOC-04 — Процедура формальной инспекции v1.0
% Комплект документации криптотренажёра mephi-12 (backend)
% Автор: Георгиев Д.Д., 24.10.2026
\documentclass[12pt,a4paper]{article}

% =============================================================================
% КОНФИГУРАЦИЯ И ПРЕАМБУЛА (единая для комплекта)
% =============================================================================

% ---------- Метаданные ----------
\newcommand{\Company}{Команда mephi-12}
\newcommand{\DocumentTitleMain}{Процедура формальной инспекции}
\newcommand{\DocumentTitleSub}{для криптотренажёра mephi-12}
\newcommand{\Version}{1.0}
\newcommand{\DocDate}{24.10.2026}
\newcommand{\AuthorName}{Георгиев Д.Д.}
\newcommand{\AgreedBy}{}
\newcommand{\ApprovedBy}{}
\newcommand{\City}{Москва}

% --- пакеты и глобальные установки ---
\usepackage[utf8]{inputenc}
\usepackage[T2A]{fontenc}
\usepackage[russian]{babel}

% Times New Roman-like
\usepackage{newtxtext,newtxmath}
\usepackage{setspace}
\onehalfspacing % межстрочный интервал 1.5

% Страницы
\usepackage{geometry}
\geometry{left=25mm,right=15mm,top=20mm,bottom=20mm}

% отступ первой строки
\usepackage{indentfirst}
\setlength{\parindent}{1.25cm}

% таблицы, картинки, списки, подписи
\usepackage{graphicx}
\usepackage{caption}
\usepackage{subcaption}
\captionsetup[table]{position=top}    % подпись таблицы над ней
\captionsetup[figure]{position=bottom}% подпись рисунка под ней
\usepackage{booktabs}
\usepackage{tabularx}
\usepackage{longtable}
\usepackage{array}
\usepackage{enumitem}
\setlist[itemize]{left=0pt,label=\textbullet}
\setlist[enumerate]{left=0pt,label=\arabic*.}

% гиперссылки
\usepackage{hyperref}
\hypersetup{colorlinks=true,linkcolor=black,urlcolor=blue}

\pagestyle{plain}  % только номер страницы внизу по центру

% нумерация разделов — стандартная десятичная (article по умолчанию)
\usepackage{tocloft} % управление оглавлением
\setlength{\cftbeforesecskip}{6pt}

% небольшие утилиты
\newcommand{\reqfield}[1]{\textbf{#1}:} % для шаблонов в приложениях

\begin{document}

% =============================================================================
% ТИТУЛЬНЫЙ ЛИСТ
% =============================================================================
\begin{titlepage}
  \thispagestyle{empty}
  \centering
  \vspace*{5.8cm}
  {\Huge \textbf{\DocumentTitleMain}\par}
  {\Large \vspace{4pt} \DocumentTitleSub \par}
  \vspace{0.8cm}

  {\large Версия: \Version \par}
  {\large Дата: \DocDate \par}
  \vspace{1.2cm}

  % ----- Таблица с ролями -----
  \centering
  \begin{tabularx}{0.95\textwidth}{@{\hspace{6pt}}p{5.5cm} p{6cm} p{4cm}@{\hspace{6pt}}}
    \toprule
    \textbf{Роль} & \textbf{ФИО} & \textbf{Подпись / Дата} \\
    \midrule
    Разработал: & \AuthorName & \\
    \addlinespace[6pt]
    Согласовано: & \AgreedBy & \\
    \addlinespace[6pt]
    Утверждено: & \ApprovedBy & \\
    \bottomrule
  \end{tabularx}

  \vfill
  {\large г. \City, \the\year \par}
\end{titlepage}

% =============================================================================
% ЛИСТ РЕГИСТРАЦИИ ИЗМЕНЕНИЙ
% =============================================================================
\section*{Лист регистрации изменений \texttt{[FI.CHANGES]}}
\addcontentsline{toc}{section}{Лист регистрации изменений}

\begin{table}[h!]
\centering
\small
\begin{tabularx}{\textwidth}{
  >{\centering\arraybackslash}p{1.1cm}
  >{\centering\arraybackslash}p{2.1cm}
  p{2.3cm}
  p{2.6cm}
  p{1.6cm}
  X
  p{1.8cm}
}
\toprule
\textbf{Версия} & \textbf{Дата} & \textbf{Автор} & \textbf{Изменённый элемент} & \textbf{Х-р изменения} & \textbf{Описание изменения} & \textbf{Утвердил} \\
\midrule
1.0 & 24.10.2026 & Георгиев Д.Д. & Весь документ & Добавлено & Первичный выпуск документа. & \\
\bottomrule
\end{tabularx}
\end{table}

\clearpage
\tableofcontents
\clearpage

% =============================================================================
% ОСНОВНОЕ СОДЕРЖАНИЕ
% =============================================================================

\section{Введение \texttt{[FI.INTRO]}}

\subsection{Назначение}

Настоящий документ (DOC-04) устанавливает обязательную процедуру формальной
инспекции (ФИ) проектной документации, разрабатываемой для backend-части
платформы «Платформа для изучения и практики курса по криптографии в МИФИ»
(краткое наименование --- криптотренажёр mephi-12). Формальная инспекция является
основным методом независимой верификации: каждый документ приоритета К
(критический путь) и В (высокий) обязан пройти ФИ до выпуска своей базовой
версии v1.0 и слияния в основную ветвь репозитория.

Документ определяет роли участников и правила их ротации, этапы процедуры,
чек-листы проверки, классификацию замечаний, метрики качества инспекции и форму
протокола. Результаты каждой инспекции фиксируются в протоколах
DOC-13.1--DOC-13.6 и в реестре замечаний \texttt{defects.md}; положительное
решение инспекции --- необходимое условие признания документа завершённым
(Definition of Done по «Процессам и обязанностям»).

\subsection{Область применения}

Процедура распространяется на все документы приоритетов К и В по Реестру
документов комплекта, включённые в расписание инспекций (раздел~\ref{subsec:schedule}):
DOC-02 (ТЗ), DOC-03 (ПУКПО), DOC-05 (системные требования), DOC-06 и
DOC-06.1--06.6 (программные требования), DOC-07 (архитектура), DOC-08 (API),
DOC-09 (матрица трассировки), DOC-10 (план верификации ПО), DOC-11 (план
интеграционного тестирования).

Объектом инспекции являются \textbf{документы} (текст, таблицы, трассировочные
матрицы, формы), а не программный код: решением команды от 20.10 работа над
проектом полностью документационная --- тестовый код не пишется и прогоны не
выполняются. Верификационные кейсы TC-XXX проверяются инспекцией как текст
планов DOC-10 и DOC-11, включая описанные на бумаге негативные проверки
(401 без JWT, 403 без роли ADMIN, 400 при невалидном теле запроса); их
исполнение не планируется, протоколы прогонов (TPR) в комплект не входят.

На документы приоритета О (DOC-14) процедура может применяться добровольно.
Сам настоящий документ (DOC-04) является приоритета К и применяется ко всем
последующим своим версиям.

\subsection{Термины и сокращения}

\begin{itemize}
  \item \textbf{Формальная инспекция (ФИ)} --- структурированная коллегиальная
        проверка артефакта по правилам настоящей процедуры (метод Фагана,
        адаптированный под учебный проект).
  \item \textbf{Артефакт} --- объект инспекции; в рамках настоящего проекта ---
        документ комплекта (исходный файл \texttt{.tex} и собранный PDF,
        CSV-источник и PDF матрицы трассировки).
  \item \textbf{Замечание (дефект), DEF} --- выявленное на инспекции
        несоответствие; регистрируется с уникальным идентификатором DEF-XXX
        (раздел~\ref{sec:def}).
  \item \textbf{Автор} --- ответственный за документ (владелец по Реестру либо
        редактор сводного документа); готовит документ к инспекции и закрывает
        замечания.
  \item \textbf{Ведущий} --- участник, который организует инспекцию, ведёт
        совещание и контролирует закрытие замечаний; не является автором
        инспектируемого документа.
  \item \textbf{Секретарь} --- участник, который протоколирует инспекцию,
        присваивает замечаниям идентификаторы DEF-XXX и оформляет протокол;
        не является автором инспектируемого документа.
  \item \textbf{Рецензент} --- участник, проверяющий документ по чек-листам
        (раздел~\ref{sec:checklists}); проверяет разделы только чужих доменов
        владения.
  \item \textbf{Домен владения} --- закреплённая за участником область кода и
        требований (карточка домена: свой раздел СТ, свои SWR, свои эндпоинты,
        свои TC).
  \item \textbf{Базовая версия} --- версия v1.0 документа, прошедшего ФИ с
        решением «ПРИНЯТО» или «ПРИНЯТО С ЗАМЕЧАНИЯМИ», слитая в \texttt{main}
        и помеченная тегом версии.
  \item \textbf{Major / Minor} --- категории серьёзности замечания
        (раздел~\ref{subsec:severity}).
  \item \textbf{SR / SWR / TC} --- системное требование (DOC-05), требование к
        ПО (DOC-06), верификационный кейс (DOC-10).
\end{itemize}

\subsection{Нормативная основа}

Процедура основана на методе инспекций М.~Фагана:
\begin{itemize}
  \item Fagan M.E. Design and Code Inspections to Reduce Software Development
        Errors // IBM Systems Journal. 1976. Vol.~15, No.~3. P.~182--211.
\end{itemize}

Из метода Фагана заимствованы: шесть этапов (планирование, обзор, подготовка,
инспекция, доработка, проверка закрытия), разделение ролей (автор, ведущий,
секретарь, рецензенты), запрет автору оценивать собственный текст, ведение
статистики замечаний и ограничение длительности совещания. Стиль требований и
роль независимой верификации соответствуют практикам DO-178C/КТ-178C
(цели верификации и обеспечения качества процессов).

Адаптация под учебный проект: инспектируются только документы (не код);
команда из шести человек с полной ротацией ролей; сокращённые нормы времени
(обзор 30--60 минут, совещание не более 2 часов, подготовка порядка 1 часа на
10 страниц). Процедура согласована с ПУКПО (DOC-03), Реестром документов и
разделом~2 «Процессы и обязанности» (расписание ФИ, git-процесс, реестр DEF).

\section{Организация инспекции \texttt{[FI.ORG]}}

\subsection{Роли и ответственности}

В каждой инспекции участвуют: автор документа, ведущий, секретарь и не менее
двух рецензентов. Обязанности и ограничения ролей приведены в
таблице~\ref{tab:roles}.

\begin{table}[h!]
  \centering
  \caption{Роли участников формальной инспекции}
  \label{tab:roles}
  \begin{tabularx}{\textwidth}{
    >{\raggedright\arraybackslash}p{2.6cm}
    >{\raggedright\arraybackslash}X
    >{\raggedright\arraybackslash}p{4.4cm}}
    \toprule
    \textbf{Роль} & \textbf{Обязанности} & \textbf{Ограничения} \\
    \midrule
    Автор &
    Готовит документ к инспекции; представляет материал на обзоре; отвечает на
    вопросы рецензентов; закрывает принятые DEF; выпускает базовую версию по
    git-процессу. &
    Не ведущий и не секретарь своего документа; не оценивает собственный текст
    как рецензент. \\
    \addlinespace
    Ведущий &
    Планирует инспекцию; проверяет критерии входа; распределяет домены
    проверки; ведёт совещание и следит за регламентом; подтверждает закрытие
    DEF; формулирует решение. &
    Не автор инспектируемого документа. \\
    \addlinespace
    Секретарь &
    Протоколирует совещание; присваивает и регистрирует DEF-XXX; ведёт
    статистику замечаний; оформляет и рассылает протокол DOC-13.X; вносит
    итоги в \texttt{defects.md}. &
    Не автор инспектируемого документа. \\
    \addlinespace
    Рецензенты (не менее 2) &
    Готовятся индивидуально по чек-листам; на совещании предъявляют и
    обосновывают замечания; проверяют факт устранения в своей части при
    повторной инспекции. &
    Проверяют разделы только \emph{чужих} доменов владения. \\
    \bottomrule
  \end{tabularx}
\end{table}

\subsection{Правило перекрёстности и ротации}

Для сводных документов (DOC-05 СТ, DOC-06 SWR), разделы которых пишут разные
владельцы доменов, действует \textbf{правило перекрёстности}: рецензент
проверяет разделы только тех доменов, которыми не владеет сам. Тем самым
исключается конфликт интересов «автор проверяет сам себя» и обеспечивается
свежий взгляд на содержание. Закрепление рецензентов за чужими доменами
выполняет ведущий на этапе планирования; для ФИ-SWR оно указано прямо в
расписании (таблица~\ref{tab:fi-schedule}).

Ротация ролей по комплекту документов: каждый из шести участников за комплект
ровно один раз выступает ведущим и ровно один раз --- секретарём; автором
инспекции является ответственный за документ (редактор), который не ведёт эту
инспекцию и не пишет её протокол.

\subsection{Расписание формальных инспекций проекта}
\label{subsec:schedule}

Расписание инспекций комплекта приведено в таблице~\ref{tab:fi-schedule}
(даты --- ноябрь 2026 г.; протокол оформляется секретарём не позднее следующего
рабочего дня после сессии).

{\small
\setlength{\tabcolsep}{3pt}
\begin{longtable}{
  >{\raggedright\arraybackslash}p{1.4cm}
  >{\raggedright\arraybackslash}p{1.8cm}
  >{\raggedright\arraybackslash}p{1.5cm}
  >{\raggedright\arraybackslash}p{1.6cm}
  >{\raggedright\arraybackslash}p{2.3cm}
  >{\raggedright\arraybackslash}p{1.5cm}
  >{\raggedright\arraybackslash}p{1.5cm}
  >{\raggedright\arraybackslash}p{3.5cm}}
\caption{Расписание формальных инспекций проекта mephi-12 (ноябрь 2026)}
\label{tab:fi-schedule}\\
\toprule
\textbf{ФИ} & \textbf{Объект} & \textbf{Сессия} & \textbf{Протокол} &
\textbf{Автор} & \textbf{Ведущий} & \textbf{Секретарь} & \textbf{Рецензенты} \\
\midrule
\endfirsthead
\toprule
\textbf{ФИ} & \textbf{Объект} & \textbf{Сессия} & \textbf{Протокол} &
\textbf{Автор} & \textbf{Ведущий} & \textbf{Секретарь} & \textbf{Рецензенты} \\
\midrule
\endhead
\bottomrule
\endfoot
\endlastfoot
ФИ-ТЗ & DOC-02 & 31.10--01.11 & 02.11, DOC-13.1 & Кокорев & Анохин &
Георгиев & Портнов, Самбьюи, Минасов \\
\addlinespace
ФИ-ПУКПО & DOC-03 & 03--04.11 & 04.11, DOC-13.2 & Самбьюи & Кокорев &
Минасов & Анохин, Георгиев \\
\addlinespace
ФИ-СТ & DOC-05 & 06.11 & 07.11, DOC-13.3 & Анохин (ред.) + авторы разделов &
Минасов & Портнов & Георгиев, Самбьюи, Кокорев \\
\addlinespace
ФИ-SWR & DOC-06, DOC-06.1--06.6 & 13--14.11 & 14.11, DOC-13.4 &
Анохин (ред.) + владельцы & Самбьюи & Кокорев &
Георгиев (чужие: аутентификация, Эль-Гамаль); Минасов (чужие: сессии, ранец);
Портнов (чужие: конфигурации, бэкдор) \\
\addlinespace
ФИ-АРХ+API & DOC-07, DOC-08 & 16--17.11 & 17.11, DOC-13.5 &
Георгиев, Минасов & Портнов & Анохин & Кокорев, Самбьюи \\
\addlinespace
ФИ-TRACE+VER & DOC-09, DOC-10, DOC-11 & 20.11 & 21.11, DOC-13.6 &
Портнов, Кокорев, Минасов & Георгиев & Самбьюи & Анохин \\
\end{longtable}
}

При изменении сроков спринтов расписание корректиется ведущим соответствующей
ФИ по согласованию с автором; изменения фиксируются в листе регистрации
изменений настоящего документа.

\section{Этапы процедуры \texttt{[FI.STAGES]}}

Процедура включает шесть этапов. Для каждого этапа определены входные условия,
действия и выход (результат). Переход к следующему этапу допускается только при
выполнении выходных критериев предыдущего; возврат на более ранний этап
допускается в любой момент по решению ведущего.

\subsection{Этап 1. Планирование}

\reqfield{Вход} документ достиг статуса «На проверке» по ПУКПО (DOC-03).

\reqfield{Действия} Ведущий совместно с автором:
\begin{itemize}
  \item назначают дату, время и место совещания, определяют состав участников
        согласно разделу~2 (автор, ведущий, секретарь, не менее двух
        рецензентов);
  \item проверяют критерии входа: (а)~документ в статусе «На проверке»;
        (б)~версия стабильна --- зафиксирована коммитом/тегом и не изменяется
        до завершения инспекции; (в)~документ собран в PDF и доступен
        участникам; (г)~версия и обозначение совпадают с Реестром документов;
  \item рассылают рецензентам материалы (исходный текст и PDF) не менее чем за
        \textbf{2 рабочих дня} до совещания; распределяют рецензентам чужие
        домены проверки.
\end{itemize}

\reqfield{Выход} назначенная инспекция в расписании команды, распределённые
роли и домены, материалы, полученные всеми рецензентами в срок. Если материалы
не разосланы за 2 рабочих дня, совещание переносится.

\subsection{Этап 2. Обзор}

\reqfield{Вход} материалы инспекции получены участниками; состав согласован.

\reqfield{Действия} Автор представляет документ рецензентам (30--60 минут):
назначение и место документа в комплекте, структуру, методику идентификаторов,
словарь терминов, известные спорные места. Ведущий раздаёт чек-листы
(раздел~\ref{sec:checklists}) и закрепляет домены проверки. Вопросы общего
контекста задаются и снимаются на этом этапе, а не на совещании.

\reqfield{Выход} рецензенты понимают контекст документа и готовы к
индивидуальной подготовке. Если на обзоре выявлены концептуальные пробелы
(документ не соответствует своему назначению по Реестру), ведущий возвращает
документ автору без выдачи DEF --- критерий входа не выполнен, инспекция
назначается заново.

\subsection{Этап 3. Подготовка}

\reqfield{Вход} обзор проведён; чек-листы и материалы у рецензентов.

\reqfield{Действия} Каждый рецензент \emph{индивидуально} читает закреплённые
за ним разделы по чек-листам раздела~\ref{sec:checklists}. Норма времени ---
порядка \textbf{1 часа на 10 страниц}. Рецензент заранее фиксирует каждое
предварительное замечание с указанием раздела/пункта документа и предлагаемой
категорией (серьёзность и тип) и направляет список секретарю до начала
совещания. Обсуждение между рецензентами на этом этапе не допускается.

\reqfield{Выход} предварительные списки замечаний у секретаря; заполненные
колонки «Да/Нет» чек-листов у рецензентов. Рецензент, не выполнивший
подготовку, к совещанию не допускается (совещание переносится).

\subsection{Этап 4. Совещание по инспекции}

\reqfield{Вход} все рецензенты подготовлены; предварительные замечания у
секретаря.

\reqfield{Действия} Ведущий ведёт проверку \emph{построчно/по разделам}, а не
«в целом». Порядок: замечания по текущему разделу предъявляются, обсуждаются и
классифицируются; секретарь фиксирует каждое замечание в протоколе с категорией
(серьёзность $\times$ тип, раздел~\ref{sec:def}) и присваивает идентификатор
DEF-XXX. По каждому замечанию совещание принимает решение:
\emph{принять} (замечание становится DEF), \emph{отклонить} (с обоснованием в
протоколе) или \emph{уточнить} (дообоснование рецензентом до конца дня).
Спорные вопросы решает ведущий. Цель совещания --- обнаружение дефектов, а не
их исправление: обсуждение способа устранения ограничивается 2--3 минутами на
замечание. Лимит совещания --- \textbf{2 часа}; непросмотренные разделы
выносятся на продолжение в другой день в пределах сессии инспекции.

\reqfield{Выход} полный перечень DEF с категориями и решениями; статистика
замечаний по категориям; черновик протокола.

\subsection{Этап 5. Доработка}

\reqfield{Вход} перечень DEF передан автору; протокол разослан участникам.

\reqfield{Действия} Автор устраняет принятые DEF в ветке \texttt{docs/DOC-XX}
по git-процессу и отмечает каждое закрытие в реестре \texttt{defects.md} с
датой. Повторную проверку выполняет ведущий: закрытие Major проверяется
полностью, Minor --- полностью или выборочно по решению ведущего. Повторное
совещание для Minor не собирается.

\reqfield{Выход} доработанная версия документа; обновлённый реестр DEF;
подтверждения ведущего по каждому закрытию.

\subsection{Этап 6. Проверка закрытия и решение}

\reqfield{Вход} автор заявил о закрытии всех Major (и, по возможности, Minor).

\reqfield{Действия} Ведущий сверяет состояние всех DEF по реестру и принимает
решение:
\begin{itemize}
  \item \textbf{ПРИНЯТО} --- открытых замечаний нет: все DEF закрыты и
        подтверждены;
  \item \textbf{ПРИНЯТО С ЗАМЕЧАНИЯМИ} --- открыты только Minor с назначенным
        сроком закрытия до конца следующего спринта; повторная ФИ не
        требуется, контроль закрытия ведёт ведущий;
  \item \textbf{НА ДОРАБОТКУ} --- есть незакрытые Major либо доработка
        привела к существенным структурным изменениям; назначается повторная
        ФИ, начиная с этапа~1 (с тем же или обновлённым составом).
\end{itemize}
При решении «ПРИНЯТО» или «ПРИНЯТО С ЗАМЕЧАНИЯМИ» документ получает статус
базовой версии: слияние в \texttt{main}, тег вида \texttt{docs/DOC-05-v1.0},
запись в листе регистрации изменений.

\reqfield{Выход} решение, зафиксированное секретарём в протоколе DOC-13.X
(форма --- приложение~А), подписи участников; реестр DEF актуализирован.

\section{Чек-листы \texttt{[FI.CHECKLISTS]}}
\label{sec:checklists}

Чек-листы применяются рецензентами на этапе подготовки (3.3) и ведущим на
совещании (3.4). Каждый пункт оценивается ответом «Да»/«Нет»; каждое «Нет»
должно сопровождаться замечанием с идентификатором DEF-XXX. Рецензент заполняет
только пункты, относящиеся к закреплённым за ним чужим доменам; остальные
пункты помечаются «Н/П» (не проверял). Заполненные чек-листы прикладываются к
протоколу инспекции.

\subsection{Чек-лист корректности}

Назначение: проверка внутренней согласованности и корректности документа.

\begin{longtable}{|p{0.8cm}|p{6.3cm}|p{1.8cm}|p{5.2cm}|}
\caption{Чек-лист корректности}\\
\hline
\textbf{№} & \textbf{Критерий проверки} & \textbf{Да/Нет} & \textbf{Замечание (DEF)} \\
\hline
\endfirsthead
\hline
\textbf{№} & \textbf{Критерий проверки} & \textbf{Да/Нет} & \textbf{Замечание (DEF)} \\
\hline
\endhead
\hline
\endfoot
\hline
\endlastfoot
1 & Документ внутренне непротиворечив: положения разных разделов не конфликтуют между собой? & & \\
\hline
2 & Каждое положение сформулировано однозначно и допускает единственное толкование? & & \\
\hline
3 & Термины и сокращения употреблены в соответствии с глоссарием комплекта; новым терминам даны определения? & & \\
\hline
4 & Все ссылки на идентификаторы (DOC-XX, SR-XX, SWR-XX, TC-XXX, DEF-XXX) корректны: объект существует и содержание ему соответствует? & & \\
\hline
5 & Числовые данные, факты и криптографические параметры (модули, размеры ключей, коды ответов HTTP, диапазоны) верны? & & \\
\hline
6 & Исходный текст LaTeX соответствует единому шаблону комплекта (преамбула, титульный лист, структура разделов)? & & \\
\hline
7 & Титульный лист, лист регистрации изменений и оглавление присутствуют и заполнены корректно? & & \\
\hline
8 & Обозначение и версия документа на титуле совпадают с Реестром документов и тегом в репозитории? & & \\
\hline
9 & Нумерация разделов, таблиц и рисунков сквозная; все перекрёстные ссылки разрешаются? & & \\
\hline
10 & Каждое требование сформулировано проверяемо: заданы измеримые критерии, а не оценочные суждения («быстро», «удобно»)? & & \\
\hline
11 & Документ не противоречит смежным документам комплекта (ТЗ, СТ, SWR, планы верификации)? & & \\
\hline
12 & Язык и стиль соответствуют нормам; опечатки не искажают смысл? & & \\
\hline
\end{longtable}

\subsection{Чек-лист полноты}

Назначение: проверка наличия всего обязательного содержания документа.

\begin{longtable}{|p{0.8cm}|p{6.3cm}|p{1.8cm}|p{5.2cm}|}
\caption{Чек-лист полноты}\\
\hline
\textbf{№} & \textbf{Критерий проверки} & \textbf{Да/Нет} & \textbf{Замечание (DEF)} \\
\hline
\endfirsthead
\hline
\textbf{№} & \textbf{Критерий проверки} & \textbf{Да/Нет} & \textbf{Замечание (DEF)} \\
\hline
\endhead
\hline
\endfoot
\hline
\endlastfoot
1 & Присутствуют и заполнены по существу все разделы, предусмотренные шаблоном и Реестром документов для документа данного типа? & & \\
\hline
2 & В версии, предъявленной на инспекцию, нет маркеров TBD, заглушек и пустых разделов? & & \\
\hline
3 & Все требования закреплённого домена (по карточке домена) отражены в документе? & & \\
\hline
4 & Для каждой функции описано требуемое поведение в терминах «что делает система», без подмены требования решением реализации? & & \\
\hline
5 & Учтены отрицательные (негативные) случаи: доступ без аутентификации (401), отсутствие требуемой роли (403), невалидное тело запроса (400)? & & \\
\hline
6 & Для каждого требования определены границы, допущения и ограничения? & & \\
\hline
7 & Разделы терминов, сокращений и использованных источников полны? & & \\
\hline
8 & Приложения, предусмотренные Реестром для данного документа (модель данных, матрицы, формы), присутствуют? & & \\
\hline
9 & Лист регистрации изменений отражает все версии, предусмотренные планом? & & \\
\hline
10 & Объём и глубина разделов достаточны для понимания документа без обращения к автору? & & \\
\hline
\end{longtable}

\subsection{Чек-лист трассируемости}

Назначение: проверка целостности цепочки SR $\to$ SWR $\to$ TC по матрице
трассировки DOC-09 и документам требований.

\begin{longtable}{|p{0.8cm}|p{6.3cm}|p{1.8cm}|p{5.2cm}|}
\caption{Чек-лист трассируемости}\\
\hline
\textbf{№} & \textbf{Критерий проверки} & \textbf{Да/Нет} & \textbf{Замечание (DEF)} \\
\hline
\endfirsthead
\hline
\textbf{№} & \textbf{Критерий проверки} & \textbf{Да/Нет} & \textbf{Замечание (DEF)} \\
\hline
\endhead
\hline
\endfoot
\hline
\endlastfoot
1 & Каждое системное требование SR трассируется не менее чем на одно программное требование SWR? & & \\
\hline
2 & Каждое системное требование SR трассируется не менее чем на один верификационный кейс TC из DOC-10? & & \\
\hline
3 & Каждое программное требование SWR трассируется не менее чем на одно SR (нет «сиротских» требований к ПО)? & & \\
\hline
4 & Каждый верификационный кейс TC привязан к требованию (нет «висящих» кейсов без требования)? & & \\
\hline
5 & Идентификаторы в матрице трассировки DOC-09 совпадают с идентификаторами в самих документах СТ (DOC-05) и SWR (DOC-06)? & & \\
\hline
6 & При переходе SR $\to$ SWR смысл требования сохранён и не сужен? & & \\
\hline
7 & Изменения требований относительно предыдущей версии отражены в матрице трассировки? & & \\
\hline
8 & Приоритеты и статусы требований согласованы между DOC-05, DOC-06 и DOC-09? & & \\
\hline
9 & CSV-источник матрицы трассировки и собранная PDF-версия синхронны по составу записей? & & \\
\hline
\end{longtable}

\section{Классификация замечаний \texttt{[FI.DEF]}}
\label{sec:def}

Каждое принятое замечание характеризуется парой «серьёзность -- тип».
Классификацию предлагает рецензент, утверждает совещание (при разногласии ---
ведущий); фиксирует секретарь.

\subsection{Серьёзность}
\label{subsec:severity}

\begin{table}[h!]
  \centering
  \caption{Категории серьёзности замечаний}
  \label{tab:severity}
  \begin{tabularx}{\textwidth}{
    >{\raggedright\arraybackslash}p{1.8cm}
    >{\raggedright\arraybackslash}X
    >{\raggedright\arraybackslash}p{4.8cm}}
    \toprule
    \textbf{Категория} & \textbf{Критерии} & \textbf{Влияние на решение} \\
    \midrule
    Major &
    Замечание искажает смысл требования; пропуск обязательного содержания;
    нарушение трассировки (разрыв или ложная связь в цепочке SR $\to$ SWR
    $\to$ TC); прямое противоречие другому документу комплекта. &
    Должно быть закрыто и подтверждено ведущим до вынесения решения; открытое
    Major влечёт «НА ДОРАБОТКУ». \\
    \addlinespace
    Minor &
    Стиль, оформление, опечатки, неидеальная, но однозначно понимаемая
    формулировка; локальные улучшения, не влияющие на смысл и трассировку. &
    Может оставаться открытым при решении «ПРИНЯТО С ЗАМЕЧАНИЯМИ» со сроком
    закрытия до конца следующего спринта. \\
    \bottomrule
  \end{tabularx}
\end{table}

\subsection{Тип замечания}

\begin{table}[h!]
  \centering
  \caption{Типы замечаний}
  \label{tab:deftype}
  \begin{tabularx}{\textwidth}{
    >{\raggedright\arraybackslash}p{2.4cm}
    >{\raggedright\arraybackslash}X}
    \toprule
    \textbf{Тип} & \textbf{Содержание} \\
    \midrule
    CONTENT & Ошибки содержания: неверные факты и выкладки, неоднозначные
    формулировки, внутренние противоречия. \\
    \addlinespace
    STD & Нарушение стандартов и оформления: отклонение от шаблона комплекта,
    обозначений, стиля, правил git-процесса. \\
    \addlinespace
    TRACE & Дефекты трассируемости: пропуски, лишние или неверные связи в
    цепочке SR $\to$ SWR $\to$ TC, рассинхрон DOC-09 с документами. \\
    \addlinespace
    COMPLETE & Неполнота: отсутствует обязательное содержание (разделы,
    требования домена, негативные случаи, приложения). \\
    \bottomrule
  \end{tabularx}
\end{table}

\subsection{Жизненный цикл замечания DEF-XXX}

Каждому принятому замечанию секретарь присваивает уникальный идентификатор
DEF-XXX со сквозной нумерацией в пределах комплекта. Отклонённые замечания
фиксируются в протоколе с обоснованием, но идентификаторов и записей в реестре
не получают. Реестр замечаний ведётся в файле \texttt{defects.md} в корне
комплекта; по каждой записи хранятся: идентификатор; документ-источник (DOC-XX)
и раздел; описание; серьёзность; тип; ответственный; статус; дата открытия;
дата закрытия; отметка о подтверждении ведущего.

\begin{table}[h!]
  \centering
  \caption{Состояния замечания в реестре}
  \label{tab:deflifecycle}
  \begin{tabularx}{\textwidth}{
    >{\raggedright\arraybackslash}p{2.4cm}
    >{\raggedright\arraybackslash}p{3.8cm}
    >{\raggedright\arraybackslash}X}
    \toprule
    \textbf{Состояние} & \textbf{Кто переводит} & \textbf{Условие перехода} \\
    \midrule
    Открыт & Секретарь, на совещании &
    Замечание принято совещанием (решение «принять» или «уточнить» с
    подтверждением); присвоен DEF-XXX, указаны документ и раздел. \\
    \addlinespace
    Назначен & Ведущий &
    Определены ответственный (автор документа) и срок закрытия. \\
    \addlinespace
    Закрыт & Ведущий &
    Исправление внесено автором и проверено ведущим; в реестре проставлены
    дата закрытия и подтверждение ведущего. \\
    \bottomrule
  \end{tabularx}
\end{table}

Срок закрытия замечаний --- \textbf{до конца следующего спринта}
(контрольные даты спринтов: 07.11, 21.11, 01.12). Открытые Minor на момент
решения «ПРИНЯТО С ЗАМЕЧАНИЯМИ» закрываются в тот же срок; контроль
осуществляет ведущий данной ФИ, отчёт о закрытии --- на очной свёрке команды в
день дедлайна спринта.

\section{Метрики и критерии качества \texttt{[FI.METRICS]}}

Метрики собирает секретарь по итогам совещания и приводит их в протоколе
DOC-13.X. Метрики служат ориентирами для ведущего при принятии решения и
для улучшения процесса; они не применяются как автоматическое правило.

\begin{table}[h!]
  \centering
  \caption{Метрики формальной инспекции}
  \label{tab:metrics}
  \begin{tabularx}{\textwidth}{
    >{\raggedright\arraybackslash}p{3.4cm}
    >{\raggedright\arraybackslash}X
    >{\raggedright\arraybackslash}p{4.6cm}}
    \toprule
    \textbf{Метрика} & \textbf{Определение} & \textbf{Ориентир (учебный проект)} \\
    \midrule
    Число замечаний по категориям &
    $N$ --- всего; разбивка по серьёзности (Major/Minor) и по типам
    (CONTENT, STD, TRACE, COMPLETE). &
    Доля Major свыше 20\% --- признак системных проблем подготовки документа. \\
    \addlinespace
    Плотность дефектов &
    $D = N/V$, где $V$ --- объём документа в страницах собранного PDF. &
    $D \le 1$ зам./стр. --- приемлемо; $D > 3$ зам./стр. --- документ предъявлен
    сырым, ведущий рассматривает «НА ДОРАБОТКУ» без полного прохода. \\
    \addlinespace
    Скорость проверки &
    $S = V/T$, где $T$ --- суммарное время подготовки и совещания, ч. &
    $S \le 10$ стр./ч; выше --- риск поверхностной проверки, повторить этап
    подготовки. \\
    \addlinespace
    Эффективность подготовки &
    $E = (N_{0}/N)\cdot100\%$, где $N_{0}$ --- число замечаний, найденных до
    совещания. &
    $E \ge 50\%$ --- подготовка работает; $E < 30\%$ --- усилить этап~3.3
    (пересмотреть чек-листы или нормы времени). \\
    \addlinespace
    Полнота закрытия &
    Доля DEF, закрытых к назначенному сроку. &
    100\% Major --- к моменту решения; Minor --- к концу следующего спринта. \\
    \bottomrule
  \end{tabularx}
\end{table}

Динамика метрик по инспекциям комплекта рассматривается командой на очных
сверках спринтов; устойчиво низкая эффективность подготовки является основанием
для корректировки настоящей процедуры (новая версия документа по git-процессу).

% =============================================================================
% ПРИЛОЖЕНИЕ
% =============================================================================
\appendix

\section{Форма протокола формальной инспекции \texttt{[FI.FORM]}}

Форма заполняется секретарём по ходу совещания, подписывается участниками и
хранится в комплекте как протокол DOC-13.X (одна инспекция --- один протокол).
Пустые поля заполняются от руки или в LaTeX; подчёркивания обозначают места
заполнения.

\subsection{Общие сведения}

\begin{table}[h!]
\centering
\begin{tabularx}{\textwidth}{|>{\raggedright\arraybackslash}p{5.2cm}|>{\raggedright\arraybackslash}X|}
\hline
\textbf{Объект инспекции (отметить)} &
ФИ-ТЗ \quad ФИ-ПУКПО \quad ФИ-СТ \quad ФИ-SWR \quad ФИ-АРХ+API \quad ФИ-TRACE+VER \\ \hline
\textbf{Инспектируемый документ (ID, наименование)} & \underline{\hspace{8cm}} \\ \hline
\textbf{Версия и дата версии} & \underline{\hspace{8cm}} \\ \hline
\textbf{Дата и время совещания} & \underline{\hspace{8cm}} \\ \hline
\textbf{Место проведения} & \underline{\hspace{8cm}} \\ \hline
\textbf{Протокол (ID)} & DOC-13.\underline{\hspace{1cm}} \\ \hline
\textbf{Применённые чек-листы (разд.~4)} & корректность \quad полнота \quad трассируемость \\ \hline
\end{tabularx}
\end{table}

\subsection{Участники и роли}

\begin{table}[h!]
\centering
\begin{tabularx}{\textwidth}{
  >{\raggedright\arraybackslash}p{3.6cm}
  >{\raggedright\arraybackslash}X
  >{\raggedright\arraybackslash}p{4.2cm}
  >{\centering\arraybackslash}p{2.4cm}}
\toprule
\textbf{Роль} & \textbf{ФИО} & \textbf{Домен проверки} & \textbf{Присутствие (Да/Нет)} \\
\midrule
Автор & & --- & \\
\addlinespace
Ведущий & & --- & \\
\addlinespace
Секретарь & & --- & \\
\addlinespace
Рецензент 1 & & & \\
\addlinespace
Рецензент 2 & & & \\
\addlinespace
Рецензент 3 & & & \\
\bottomrule
\end{tabularx}
\end{table}

\subsection{Статистика замечаний}

\begin{table}[h!]
\centering
\begin{tabularx}{\textwidth}{|>{\raggedright\arraybackslash}p{2.8cm}|>{\centering\arraybackslash}X|>{\centering\arraybackslash}X|>{\centering\arraybackslash}X|>{\centering\arraybackslash}X|>{\centering\arraybackslash}p{1.8cm}|}
\hline
\textbf{Серьёзность} & \textbf{CONTENT} & \textbf{STD} & \textbf{TRACE} & \textbf{COMPLETE} & \textbf{Всего} \\
\hline
Major & & & & & \\
\hline
Minor & & & & & \\
\hline
\textbf{Итого} & & & & & \\
\hline
\end{tabularx}
\end{table}

\subsection{Перечень замечаний}

\begin{longtable}{|
  >{\raggedright\arraybackslash}p{1.4cm}
  >{\raggedright\arraybackslash}p{1.6cm}
  >{\raggedright\arraybackslash}p{4.2cm}
  >{\raggedright\arraybackslash}p{1.5cm}
  >{\raggedright\arraybackslash}p{1.6cm}
  >{\raggedright\arraybackslash}p{1.9cm}
  >{\raggedright\arraybackslash}p{1.3cm}}
\caption{Перечень замечаний инспекции}\\
\hline
\textbf{ID} & \textbf{Раздел} & \textbf{Описание} & \textbf{Серьёзность} &
\textbf{Тип} & \textbf{Ответственный} & \textbf{Срок} \\
\hline
\endfirsthead
\hline
\textbf{ID} & \textbf{Раздел} & \textbf{Описание} & \textbf{Серьёзность} &
\textbf{Тип} & \textbf{Ответственный} & \textbf{Срок} \\
\hline
\endhead
\hline
\endfoot
\hline
\endlastfoot
DEF-\underline{\hspace{0.6cm}} & & & & & & \\
\hline
DEF-\underline{\hspace{0.6cm}} & & & & & & \\
\hline
DEF-\underline{\hspace{0.6cm}} & & & & & & \\
\hline
DEF-\underline{\hspace{0.6cm}} & & & & & & \\
\hline
DEF-\underline{\hspace{0.6cm}} & & & & & & \\
\hline
DEF-\underline{\hspace{0.6cm}} & & & & & & \\
\hline
\end{longtable}

\subsection{Решение инспекции}

По результатам инспекции принято следующее решение (отметить одно):

\medskip
\noindent( )~\textbf{ПРИНЯТО} --- документ принимается без изменений; все
замечания закрыты и подтверждены на момент решения.

\medskip
\noindent( )~\textbf{ПРИНЯТО С ЗАМЕЧАНИЯМИ} --- открыты только Minor;
срок закрытия --- до конца следующего спринта; контроль закрытия --- ведущий;
повторная инспекция не требуется.

\medskip
\noindent( )~\textbf{НА ДОРАБОТКУ} --- имеются незакрытые Major либо
существенная переработка; назначается повторная формальная инспекция.

\medskip
\noindentДата повторной инспекции (для «НА ДОРАБОТКУ»):
\underline{\hspace{4cm}} \quad Метрики: $D = \underline{\hspace{1cm}}$ зам./стр.;
$S = \underline{\hspace{1cm}}$ стр./ч; $E = \underline{\hspace{1cm}}\%$.

\subsection{Подписи участников}

\begin{table}[h!]
\centering
\begin{tabularx}{0.95\textwidth}{
  >{\raggedright\arraybackslash}p{3.2cm}
  >{\raggedright\arraybackslash}X
  >{\centering\arraybackslash}p{3.6cm}
  >{\centering\arraybackslash}p{2.4cm}}
\toprule
\textbf{Роль} & \textbf{ФИО} & \textbf{Подпись} & \textbf{Дата} \\
\midrule
Автор & & & \\
\addlinespace
Ведущий & & & \\
\addlinespace
Секретарь & & & \\
\addlinespace
Рецензент & & & \\
\addlinespace
Рецензент & & & \\
\bottomrule
\end{tabularx}
\end{table}

\end{document}
